\section{The Institution: Scope, Height, and Exit}

A consent requirement is conventionally described by two properties: which
borrowing it covers, and how large a majority it demands. Read this way, one
requirement is stronger than another when it reaches more debt or sets a higher
bar. That description is incomplete, and its missing third term is where the
institution's incidence is in fact decided. Alongside its \emph{scope} and its
\emph{height}, a consent requirement is defined by the \emph{exit menu} it
leaves open: the set of financing arrangements through which a government can
reach the same objective without facing the electorate at all. Two requirements
identical in scope and height are not the same institution if one is surrounded
by cheap exits and the other by none.

\subsection{Scope and height}

Every state draws the scope line through the form of the pledge.
General-obligation debt, secured by the government's full taxing power and
therefore falling on every taxpayer, sits inside the requirement; revenue
obligations, leases, and the debt of statutory districts and authorities,
repaid from charges on identifiable users, sit outside it. Height then fixes,
for debt inside the line, the share of the electorate that must agree: a simple
majority, fifty-five per cent, or two thirds. Scope decides whether a gate
exists; height decides how heavy it is. Neither asks what the money is for. The
requirement is triggered by how a debt will be repaid, not by the merit of the
project it finances.

\subsection{Economic feasibility: $R_j$}

Whether a good can leave the voted channel begins with a property of the good
itself. Let $R_j$ denote the economic feasibility of dedicated repayment for
activity $j$: the capacity to identify the beneficiaries of an investment and
charge them in a way that can service debt on its own. Water systems, car
parks, and power utilities score high on $R_j$; their users are metered and
billed, and the revenue they generate can stand behind a bond without recourse
to the tax base. Schools, parks, and public buildings score low; their benefits
cannot be confined to those who pay, and their debt must therefore rest on the
taxing power. $R_j$ is a characteristic of the activity, not of the government
that undertakes it, and it varies across the functions of the local state far
more than across the states themselves.

\subsection{The legal menu, $L_{is}$, and where it came from}

Economic feasibility is necessary for exit but not sufficient. A chargeable
activity escapes the ballot only if the government undertaking it holds a legal
instrument capable of exploiting the charge. Let $L_{is}$ denote the
alternative financing arrangements lawfully available to a government of type
$i$ in state $s$: the revenue obligations, lease structures, public
authorities, and special-purpose districts that the law of that state places
outside the consent requirement. $L_{is}$ is a property of law rather than of
behaviour, and it is measured here from statute and doctrine, not from the
financing a government happened to choose.

The legal menu is not a neutral financial technology, and this is the point at
which the institution's history enters its mechanism rather than merely
preceding it. The exempt instruments were built, one at a time, out of
legislation and constitutional interpretation: the special-fund doctrine that
reclassified revenue-backed obligations as something other than debt in the
constitutional sense, the lease exceptions that followed, and the statutes that
chartered authorities and districts. Contemporaries recognised the manoeuvre
early, describing public building authorities as a costly subversion of
constitutional debt limits by the late 1950s (Morris 1958). What the menu adds
to a government is therefore not only a cheaper security but, in some cases, a
different politics.

Some exits change the very franchise. After the Supreme Court struck property
qualifications in general bond elections (Cipriano v. City of Houma 1969; City
of Phoenix v. Kolodziejski 1970), it upheld property-weighted, and even
acreage-weighted, voting in specialised districts (Salyer Land Co. v. Tulare
Lake Basin Water Storage District 1973; Ball v. James 1981). A general-purpose
government that routes a project through such a district does not merely avoid a
majority of its residents; it can substitute an electorate constituted on
ownership for one constituted on residence. The freeholder franchise the
general channel lost survives, lawfully, in the exit. $L_{is}$ thus carries a
distributive history that $R_j$ does not: it records which governments were
granted an alternative, and on whose consent that alternative runs. This is why
the exit menu is a political institution rather than a financing detail, and
why the legal menu, not the economics of the good, is the term through which
the requirement's incidence acquires its political character.

\subsection{Effective exit: the match $R_j \times L_{is}$}

Exit requires both terms at once. An economically chargeable project provides
no outside option to a government that lacks a legal instrument to carry it, and
a rich legal menu is of no use to a project that cannot generate the repayment
stream the instrument needs. The effective outside option is therefore the
match between the two,

\[
E_{ijs} = g(R_j, L_{is}),
\qquad
\frac{\partial^2 E_{ijs}}{\partial R_j \, \partial L_{is}} > 0,
\]

with economic feasibility and legal availability entering as complements. In
the simplest representation, $E_{ijs} = R_j \times L_{is}$, so that exit is
available only where both are present. A school district, whose output is
almost entirely non-chargeable, has a low $R_j$ whatever its legal menu; a city
rich in chargeable services still cannot route a schoolhouse off the ballot.
The complementarity is what makes the requirement's reach uneven in a way that
neither the economics nor the law would predict alone, and it is what accounts
for the cases a purely functional account leaves awkward: townships that borrow
like small school districts, or general-purpose governments whose few
non-chargeable projects remain exposed while the rest of their capital escapes.

The institution is therefore properly characterised by three terms rather than
two,

\[
\text{Institution} = (\text{Scope}, \text{Height}, \text{Exit Menu}),
\]

and the response a government makes to a consent requirement is a function not
of the rule's stringency alone but of the requirement it faces and the exit
available to it,

\[
\text{Response} = f(\theta_s, E_{ijs}),
\]

where $\theta_s$ is the formal stringency of the requirement in state $s$ and
$E_{ijs}$ the quality of the outside option. The central implication follows
directly and organises the rest of the paper. The incidence of the requirement
cannot be read from its scope and height, because two governments facing the
identical rule confront it across entirely different choice sets. Where
$E_{ijs}$ is high the rule governs one instrument among many; where it is near
zero the rule is the gatekeeper of provision.
